% Draft appendix section for NeurIPS paper — node types + subclass.
% Copy to Manuscript/authomathics.tex appendix when Drive is accessible.

\subsection{Node Type Taxonomy}
\label{app:node-types}

Each node in the knowledge graph represents one mathematical object, typed
into one of 14 categories organized in four groups. Constants and groupings
are defined in \texttt{src/node\_types.py}.

\paragraph{Theorem-like nodes} state a mathematical claim with separable
assumptions and conclusions: \texttt{theorem}, \texttt{lemma},
\texttt{proposition}, \texttt{corollary}, \texttt{claim}, \texttt{conjecture}.

\paragraph{Foundational nodes:} \texttt{definition} (introduces a named
concept) and \texttt{proof} (establishes a theorem-like node, linked via
\texttt{proves} edge).

\paragraph{Pedagogical nodes} (primarily in textbooks): \texttt{example}
(worked illustration, linked via \texttt{illustrates}) and \texttt{exercise}
(reader problem, linked via \texttt{tests}).

\paragraph{Other content:} \texttt{problem} (named mathematical problem or
optimization formulation, linked to algorithms via \texttt{solves}),
\texttt{equation} (labeled/numbered), \texttt{algorithm} (pseudocode or
procedure), \texttt{remark} (author observation).

\paragraph{Subclass field.}
Every node carries an optional \texttt{subclass} field giving a finer
categorization within the type. Selected examples:

\begin{center}
\small
\begin{tabular}{llp{7cm}}
\toprule
\textbf{Type} & \textbf{Subclass} & \textbf{Description / Example} \\
\midrule
definition & \texttt{concept} & Standard math definition (vector space, group) \\
definition & \texttt{model} & Statistical / probabilistic model (MoG, LDA, HMM) \\
definition & \texttt{estimator} & Statistical estimator formula (MLE, MAP, LASSO) \\
definition & \texttt{loss\_function} & Objective function (cross-entropy, ELBO) \\
\midrule
theorem & \texttt{existence} & ``There exists \ldots'' \\
theorem & \texttt{bound} & ``$\|x\| \leq C$'' \\
theorem & \texttt{convergence} & ``$\hat\theta_n \to \theta^*$'' \\
theorem & \texttt{impossibility} & ``No algorithm can \ldots'' \\
\midrule
proof & \texttt{contradiction} & Proof by contradiction \\
proof & \texttt{induction} & Proof by induction \\
proof & \texttt{probabilistic} & Probabilistic method \\
\midrule
algorithm & \texttt{exact} & Exact solver \\
algorithm & \texttt{approximation} & Approximation algorithm (with ratio) \\
algorithm & \texttt{heuristic} & No worst-case guarantee \\
\midrule
problem & \texttt{optimization} & ``Find $x$ maximizing $g(x)$ s.t.\ $Ax \leq b$'' \\
problem & \texttt{decision} & ``Does there exist \ldots?'' (yes/no) \\
\bottomrule
\end{tabular}
\end{center}

This field is particularly important for statistics and ML papers that
present mathematical content outside the classical theorem--proof structure.
A \emph{statistical model} (e.g., a mixture of Gaussians or latent Dirichlet
allocation) is a \texttt{definition} with \texttt{subclass=model}; the MLE
is a \texttt{definition} with \texttt{subclass=estimator}; the EM algorithm
is an \texttt{algorithm} with \texttt{subclass=iterative} that
\texttt{solves} the estimation \texttt{problem}. These annotations enable
fine-grained cross-paper matching: two papers defining the same generative
model under different names (e.g., LDA vs.\ STRUCTURE) are both
\texttt{definition/model} nodes.

\paragraph{Distinction: \texttt{exercise} vs.\ \texttt{problem}.}
An \texttt{exercise} is pedagogical (the author knows the answer). A
\texttt{problem} is a mathematical object being studied (may be open, may
be an optimization formulation). Exercises connect via \texttt{tests} edges;
problems connect via \texttt{solves} edges from algorithms.
